% =====================================================================
%  REGISTRO DELLE MODIFICHE di glossario.tex
%
%  Le righe le aggiunge il bot dopo il merge di una pull request con
%  un'etichetta di versione (correzione, contenuto, approvazione):
%  GUIDA.md, sezione 5. La riga più recente è in cima.
%  Si può correggere a mano (per esempio un refuso in una descrizione)
%  con una pull request SENZA etichetta. La classe legge questo file da
%  sola: nel .tex non serve scrivere nulla.
% =====================================================================
\modifica{1.0.0}{2026-10-29}{spagnoli}{-}{Approvazione del documento}
\modifica{0.5.0}{2026-10-29}{guglielmo}{cescobolla}{Corretti la voce Stack tecnologico, l'ordine alfabetico della lettera S e le iniziali maiuscole delle voci}
\modifica{0.4.0}{2026-10-29}{biasini}{cescobolla}{Aggiunte le voci Ciclo di vita, Committente, Dominio applicativo, Fornitore, Fornitura, Proponente e Stack tecnologico; riscritta la voce Requisito}
\modifica{0.3.0}{2026-10-28}{spagnoli}{cescobolla}{Corretti i collegamenti dall'indice alle singole lettere}
\modifica{0.2.0}{2026-10-28}{spagnoli}{corro}{Aggiunte le voci Best practice, Standard, Stakeholder e WoW}
\modifica{0.1.1}{2026-10-27}{guglielmo}{corro}{Documento spostato tra i documenti interni (uso interno)}
\modifica{0.1.0}{2026-10-26}{guglielmo}{corro}{Prima stesura: voci dei termini usati nell'Analisi dei Capitolati}
